% Part: applied-modal-logics
% Chapter: epistemic-logic
% Section: public-announcement-logic
% [tr source emendation] Upstream repeats `language-epistemic-logic` here;
% this file and its content form the public-announcement-logic section.

\documentclass[../../../include/open-logic-section]{subfiles}

\begin{document}

\olfileid[tr]{aml}{el}{pal}

\olsection{Kamusal Duyuru Mantığı}

Dinamik epistemik mantıklar, öznelerin bilgisinin zaman içinde ya da yeni
bilgi edindikçe nasıl değiştiğini temsil etmemize olanak verir. Bunların
çoğu, bilgideki değişiklikleri bilgisel \emph{olaylar} ya da
\emph{güncellemeler} kullanarak temsil eder. En temel güncelleme türü, bir
formülün doğru bir biçimde kamuya duyurulduğu ve bütün öznelerin buna birlikte
tanık olduğu kamusal duyurudur. Bunu yapmak için dili şöyle genişletiriz:
% [tr source emendation] Upstream misspells `knowledge` as `knowlege` in this
% paragraph; Turkish renders the intended term.

\begin{defn}
$G$ bir özne simgeleri kümesi olsun. Kamusal duyurular içeren çok özneli
epistemik mantığın temel dili şunları içerir:
\begin{enumerate}
  \tagitem{prvFalse}{!!{falsity} için önerme sabiti~$\lfalse$.}{}
  \tagitem{prvTrue}{!!{truth} için önerme sabiti~$\ltrue$.}{}
  \item !!{denumerable}s bir !!{propositional variable}s kümesi: $\Obj
    p_0$, $\Obj p_1$, $\Obj p_2$, \dots
  \item Önerme bağlaçları: \startycommalist
  \iftag{prvNot}{\ycomma $\lnot$ (değilleme)}{}%
  \iftag{prvAnd}{\ycomma $\land$ (tümel evetleme)}{}%
  \iftag{prvOr}{\ycomma $\lor$ (tikel evetleme)}{}%
  \iftag{prvIf}{\ycomma $\lif$ (!!{conditional})}{}%
  \iftag{prvIff}{\ycomma $\liff$ (!!{biconditional})}{}.
  % [tr source emendation] Upstream omits the biconditional from this initial
  % connective inventory but admits it in the inductive grammar below and in
  % the unextended epistemic language. Restore the tagged connective here.
  \item Bilgi işleci $\Knows_a$; burada $a \in G$.
  \item Kamusal duyuru işleci $[!B]$; burada $!B$ !!a{formula}{}dür.
\end{enumerate}
\end{defn}

Kamusal duyuru işleci bir kutu işleci gibi çalışır; bu nedenle dilin
tümevarımlı tanımı şöyle verilir:

\begin{defn}
  Epistemik dilin \emph{!!^{formula}s{}i} tümevarımlı olarak şöyle tanımlanır:
  \begin{enumerate}
  \tagitem{prvFalse}{$\lfalse$ atomik bir !!{formula}{}dür.}{}

  \tagitem{prvTrue}{$\ltrue$ atomik bir !!{formula}{}dür.}{}

  \item Her önerme değişkeni $\Obj p_i$, (atomik) bir
  !!{formula}{}dür.

  \tagitem{prvNot}{$!A$ !!a{formula} ise $\lnot !A$ da
  !!a{formula}{}dür.}{}

  \tagitem{prvAnd}{$!A$ ve $!B$ !!{formula}s ise $(!A \land
  !B)$ de !!a{formula}{}dür.}{}

  \tagitem{prvOr}{$!A$ ve $!B$ !!{formula}s ise $(!A \lor !B)$
  de !!a{formula}{}dür.}{}

  \tagitem{prvIf}{$!A$ ve $!B$ !!{formula}s ise $(!A \lif !B)$
  de !!a{formula}{}dür.}{}

  \tagitem{prvIff}{$!A$ ve $!B$ !!{formula}s ise $(!A \liff !B)$
  de !!a{formula}{}dür.}{}

  \item $!A$ !!a{formula} ve $a \in G$ ise $\Knows_a !A$ da
  !!a{formula}{}dür.
  
  \item $!A$ ile $!B$ !!{formula}s ise $[!A] !B$ de !!a{formula}{}dür.

  \tagitem{limitClause}{Başka hiçbir şey !!a{formula} değildir.}{}
\end{enumerate}
\end{defn}

!!{formula} $[!A] !B$ için amaçlanan okuma, ``$!A$ doğru bir biçimde kamuya
duyurulduktan sonra $!B$ geçerlidir'' biçimindedir. Kamusal duyurular
bağlamında ortak bilgiden söz etmek de zaman zaman yararlı olacağından dil,
ortak bilgi işleci $\CKnows_G !A$'yı da içerebilir.

\end{document}
